\subsection{SEGMENTS OF A WELL-ORDERED SET}

A relation \(R\{x, y\}\) is said to be a well-ordering relation between x and y if R is an order relation between x and y and if, for each non-empty set E on which \(R\{x, y\}\) induces an order relation (i.e., such that \(x \in E\) implies \(R\{x, x\}\) ; cf.~§1, no. 1), E, ordered by this relation, has a least element.

A set E ordered by an ordering \(\Gamma\) is said to be well-ordered if the relation \(y \in \Gamma\langle x \rangle\) is a well-ordering relation between x and y; \(\Gamma\) is then said to be a well-ordering on E. The following definition is equivalent to this:

DEFINITION 1. A set E is said to be well-ordered if it is ordered and if each non-empty subset of E has a least element.

A well-ordered set E is totally ordered because every subset \(\{x, y\}\) of E has a least element. Every subset A of E which is bounded above in E has a least upper bound in E.

\minorheading{Examples}

\begin{enumerate}
\def\labelenumi{(\arabic{enumi})}
\item
  Let \(\mathbf{E} = \{\alpha, \beta\}\) be a set whose elements are distinct. It is easily verified that the subset \(\{(\alpha, \alpha), (\beta, \beta), (\alpha, \beta)\}\) of \(\mathbf{E} \times \mathbf{E}\) is the graph of a well-ordering on \(\mathbf{E}\) .
\item
  Every subset (in particular, the empty subset) of a well-ordered set is well-ordered by the induced ordering.
\end{enumerate}

* (3) The existence of totally ordered sets which are not well-ordered is equivalent to the axiom of infinity (§ 4, no. 4, Corollary 1 to Proposition 3, and Exercise 3).

\begin{enumerate}
\def\labelenumi{(\arabic{enumi})}
\setcounter{enumi}{3}
\item
  If \(\Gamma\) is a well-ordering on \(E\) , the ordering opposite to \(\Gamma\) is a well-ordering on \(E\) only if \(E\) is finite ( \(\S 4\) , Exercise 3). *
\item
  Let E be a well-ordered set. The set \(E_{1}\) obtained by adjoining to E a greatest element b (§1, no. 7) is well-ordered, for if H is any non-empty subset of \(E_{1}\) other than \(\{b\}\) , the least element of \(H \cap E\) is also the least element of H.
\end{enumerate}

Remark. * As a consequence of the axiom of infinity (§6, no. 1), there exist well-ordered sets which have no greatest element, for example the set N of natural integers. *

DEFINITION 2. In an ordered set E, a subset of E such that the relations \(x \in S\) , \(y \in E\) , and \(y \leqslant x\) imply \(y \in S\) is called a segment of E.

Clearly, every intersection and union of segments of E is a segment of E. If S is a segment of E, every segment of S is also a segment of E. The set E itself and the empty set are segments of E.

PROPOSITION 1. In a well-ordered set E, every segment of E other than E itself is an interval {]}←, a{[}, where a ∈ E.

Let S be a segment of E such that \(S \neq E\) . Since E --- S is not empty, it has a least element a. By virtue of Definition 2, the relation \(x \geqslant a\) implies \(x \notin S\) ; otherwise we would have \(a \in S\) , which is absurd. Hence E --- S is the interval \([a, \rightarrow[, \text{ and } S \text{ is the interval}] \leftarrow, a[\) .

For every element \(x\) in a totally ordered set \(\mathbf{E}\) the segment \(] \leftarrow, x[\) is called the segment with endpoint \(x\) , and is denoted by \(\mathbf{S}_x\) .

Note that if E is well-ordered and not empty, \(S_{x}\) has a least element \(\alpha\) and consequently is also the interval \([\alpha, x[\) .

Let E be a totally ordered set. The union A of the \(S_{x}\) , as x runs through E, is E if E has no greatest element; and if E has a greatest element b, we have \(A = E - \{b\}\) .

PROPOSITION 2. The set \(\mathbf{E}^*\) of segments of a well-ordered set \(\mathbf{E}\) is well-ordered by inclusion. The mapping \(x \to S_x\) is an isomorphism of the well-ordered set \(\mathbf{E}\) onto the set of segments of \(\mathbf{E}\) other than \(\mathbf{E}\) itself.

It is clear that if \(x \in \mathbf{E}\) and \(y \in \mathbf{E}\) , the relation \(x \leqslant y\) implies \(\mathrm{S}_x \subset \mathrm{S}_y\) , and that \(x < y\) implies \(\mathrm{S}_x \neq \mathrm{S}_y\) ; the mapping \(x \to \mathrm{S}_x\) is therefore an isomorphism of \(\mathbf{E}\) onto the set \(\mathrm{S}(\mathbf{E})\) of segments of \(\mathbf{E}\) distinct from \(\mathbf{E}\) itself (§ 1, no. 12, Proposition 11), and consequently \(\mathrm{S}(\mathbf{E})\) is well-ordered. Moreover, \(\mathbf{E}^*\) is isomorphic to the well-ordered set obtained from \(\mathrm{S}(\mathbf{E})\) by adjoining a greatest element.

PROPOSITION 3. Let \((\mathbf{X}_{\iota})_{\iota \in \mathbf{I}}\) be a family of well-ordered sets such that for each pair of indices \((\iota, \varkappa)\) one of the sets \(\mathbf{X}_{\iota}, \mathbf{X}_{\varkappa}\) is a segment of the other. Then there exists a unique ordering on the set \(\mathbf{E} = \bigcup_{\iota \in \mathbf{I}} \mathbf{X}_{\iota}\) which induces the given ordering on each of the \(\mathbf{X}_{\iota}\) . Endowed with this ordering, \(\mathbf{E}\) is a well-ordered set. Every segment of \(X_{t}\) is a segment of E; for each \(x \in X_{t}\) , the segment with endpoint x in \(X_{t}\) is equal to the segment with endpoint x in E; and each segment of E is either E itself or a segment of one of the \(X_{t}\) .

The first assertion is a consequence of the following general lemma :

LEMMA 1. Let \((\mathbf{X}_{\alpha})_{\alpha \in \Lambda}\) be a family of ordered sets, directed with respect to the relation \(c\) (in other words, such that for each pair of indices \(\alpha, \beta\) there exists an index \(\gamma\) such that \(\mathbf{X}_{\alpha} \subset \mathbf{X}_{\gamma}\) and \(\mathbf{X}_{\beta} \subset \mathbf{X}_{\gamma}\) ). Suppose that, for each pair of indices \((\alpha, \beta)\) such that \(\mathbf{X}_{\alpha} \subset \mathbf{X}_{\beta}\) , the ordering induced on \(\mathbf{X}_{\alpha}\) by that of \(\mathbf{X}_{\beta}\) is identical with the given ordering on \(\mathbf{X}_{\alpha}\) . Under these conditions there exists a unique ordering on the set \(\mathbf{E} = \bigcup_{\alpha \in \Lambda} \mathbf{X}_{\alpha}\) which induces the given ordering on each \(\mathbf{X}_{\alpha}\) .

Let \(G_{\alpha}\) be the graph of the given ordering on \(X_{\alpha}\) . If G is the graph of an ordering on E which induces on each \(X_{\alpha}\) the ordering whose graph is \(G_{\alpha}\) , then we must have \(G_{\alpha} \subset G\) for each \(\alpha \in A\) ; hence G contains \(\bigcup_{\alpha \in A} G_{\alpha}\) . On the other hand, for each pair \((x, y)\) of elements of E there exists by hypothesis an index \(\alpha \in A\) such that \(x \in X_{\alpha}\) and \(y \in X_{\alpha}\) ; if \((x, y) \in G\) , we have \((x, y) \in G_{\alpha}\) , so that \(G \subset \bigcup_{\alpha \in A} G_{\alpha}\) . Hence if the required ordering on E exists, its graph is necessarily \(G = \bigcup_{\alpha \in A} G_{\alpha}\) . It remains to be shown that this set satisfies the conditions of the lemma. Since \(G_{\beta} \cap (X_{\alpha} \times X_{\alpha}) = G_{\alpha}\) if \(X_{\alpha} \subset X_{\beta}\) , we have \(G \cap (X_{\alpha} \times X_{\alpha}) = G_{\alpha}\) for all \(\alpha \in A\) ; on the other hand, it follows from the hypothesis that any three elements \(x, y, z\) of E belong to the same \(X_{\alpha}\) . Hence \((x, y) \in G\) is an order relation on E, and the lemma is proved.

We now take up the proof of Proposition 3. Let us begin by showing that each \(\mathbf{X}_{\iota}\) is a segment of E. Indeed, if \(x \in \mathbf{X}_{\iota}, y \in \mathbf{E}\) , and \(y \leqslant x\) , there exists an index \(x\) such that \(\mathbf{X}_{\iota} \subset \mathbf{X}_{x}\) and \(y \in \mathbf{X}_{x}\) ; since by hypothesis \(\mathbf{X}_{\iota}\) is a segment of \(\mathbf{X}_{x}\) , we have \(y \in \mathbf{X}_{\iota}\) , which proves the assertion. The same reasoning proves that for each \(x \in \mathbf{X}_{\iota}\) the segment with endpoint \(x\) in \(\mathbf{X}_{\iota}\) is identical with the interval {]}←, \(x[\) in E. Next let us show that E is well-ordered. If H is a non-empty subset of E, there is an index \(\iota \in I\) such that \(\mathrm{H} \cap \mathrm{X}_{\iota} \neq \emptyset\) ; if \(a\) is the least element of \(\mathrm{H} \cap \mathrm{X}_{\iota}\) in \(\mathbf{X}_{\iota}\) , then \(a\) is also the least element of H in E. That is, if \(x \in H\) , there exists an index \(x \in I\) such that \(\mathbf{X}_{\iota} \subset \mathbf{X}_{x}\) and \(x \in \mathbf{X}_{x}\) ; we cannot have \(x < a\) , because the interval {]}←, \(a[\) is contained in \(\mathbf{X}_{\iota}\) , and therefore we have \(x \geqslant a\) since \(\mathbf{X}_{x}\) is totally ordered.

Finally, we must show that a segment of E, other than E itself, is a segment of one of the \(X_{t}\) ; this is an immediate consequence of the preceding arguments, since such a segment is of the form \(]\leftarrow\) , \(x[\text{(Proposition 1)}\) and since x belongs to some \(X_{t}\) .

